% Part: normal-modal-logic
% Chapter: completeness
% Section: complete-consistent-sets

\documentclass[../../../include/open-logic-section]{subfiles}

\begin{document}

\olfileid[hi]{nml}{com}{ccs}

\olsection{पूर्ण $\Sigma$-संगत समुच्चय}

मान लें कि $\Sigma$ मोडल !!{formula}ों का समुच्चय है---इन्हें किसी सामान्य
मोडल तर्कशास्त्र के अभिगृहीत या परिभाषक सिद्धांत समझें। कोई समुच्चय
$\Gamma$, $\Sigma$-संगत तभी है जब $\Gamma \Proves/[\Sigma] \lfalse$,
अर्थात् जब $\Sigma$ से $!A_1 \lif (!A_2 \lif \cdots (!A_n \lif
\lfalse)\dots)$ की कोई !!{derivation} न हो, जहाँ प्रत्येक $!A_i \in
\Gamma$। हम ऐसा ``विहित'' प्रतिरूप बनाएँगे जिसमें हर जगत एक विशेष प्रकार
का $\Sigma$-संगत समुच्चय होगा: ऐसा समुच्चय जो केवल $\Sigma$-संगत ही नहीं,
बल्कि इस अर्थ में अधिकतम भी हो कि वह प्रत्येक मोडल !!{formula} का सत्यता-मूल्य
निर्धारित करता है---हर $!A$ के लिए या तो $!A \in \Gamma$, या $\lnot !A
\in \Gamma$:

\begin{defn}
  कोई समुच्चय $\Gamma$ \emph{पूर्ण $\Sigma$-संगत} तभी है जब वह
  $\Sigma$-संगत हो और प्रत्येक~$!A$ के लिए या तो $!A \in \Gamma$, या
  $\lnot !A \in \Gamma$ हो।
\end{defn}

पूर्ण $\Sigma$-संगत समुच्चयों $\Gamma$ के अनेक उपयोगी गुण हैं। एक तो वे
निगमनतः संवृत हैं, अर्थात् यदि $\Gamma \Proves[\Sigma] !A$, तो $!A \in
\Gamma$। विशेषतः इसका अर्थ है कि किसी !!a{formula}~$!A \in \Sigma$ का हर
दृष्टांत $\Gamma$ में भी है। इसके अतिरिक्त, $\Gamma$ में सदस्यता
प्रतिज्ञप्तिक संयोजकों की सत्यता-शर्तों को प्रतिबिंबित करती है। ``विहित
प्रतिरूप'' परिभाषित करते समय यह महत्त्वपूर्ण होगा।

\begin{prop}\ollabel{prop:ccs-properties}
  मान लें कि $\Gamma$ पूर्ण $\Sigma$-संगत है। तब:
  \begin{enumerate}
  \item \ollabel{prop:ccs-closed}%
    $\Gamma$, $\Sigma$ में निगमनतः संवृत है।
  \item \ollabel{prop:ccs-sigma}%
    $\Sigma \subseteq \Gamma$.
  \tagitem{prvFalse}{\ollabel{prop:ccs-lfalse}%
    $\lfalse \notin \Gamma$}{}
  \tagitem{prvTrue}{\ollabel{prop:ccs-ltrue}%
    $\ltrue \in \Gamma$}{}
  \item \ollabel{prop:ccs-lnot}%
    $\lnot!A \in \Gamma$ तभी जब $!A \notin \Gamma$।
  \tagitem{prvAnd}{\ollabel{prop:ccs-land}%
    $!A \land !B \in \Gamma$ तभी जब $!A \in \Gamma$ और $!B \in \Gamma$}{}
  \tagitem{prvOr}{\ollabel{prop:ccs-lor}%
    $!A \lor !B \in \Gamma$ तभी जब $!A \in \Gamma$ या $!B \in \Gamma$}{}
  \tagitem{prvIf}{\ollabel{prop:ccs-lif}%
    $!A \lif !B \in \Gamma$ तभी जब $!A \notin \Gamma$ या $!B \in \Gamma$}{}
  \tagitem{prvIff}{\ollabel{prop:ccs-liff}%
    $!A \liff !B \in \Gamma$ तभी जब या तो $!A \in \Gamma$ और $!B \in
    \Gamma$, या $!A \notin \Gamma$ और $!B \notin \Gamma$}{}
  \end{enumerate}
\end{prop}

\begin{proof}
  \begin{enumerate}
  \item मान लें कि $\Gamma \Proves[\Sigma] !A$, किंतु $!A \notin
    \Gamma$। चूँकि $\Gamma$ पूर्ण $\Sigma$-संगत है, $\lnot!A \in
    \Gamma$। इससे $\Gamma$ असंगत होगा, क्योंकि $!A, \lnot !A
    \Proves[\Sigma] \lfalse$।
    
  \item यदि $!A \in \Sigma$, तो $\Gamma \Proves[\Sigma] !A$, और
    निगमनात्मक संवृत्ति, अर्थात् स्थिति~\olref{prop:ccs-closed}, से
    $!A \in \Gamma$।
    
  \tagitem{prvFalse}{यदि $\lfalse \in \Gamma$, तो $\Gamma
    \Proves[\Sigma] \lfalse$, इसलिए $\Gamma$, $\Sigma$-असंगत होगा।}{}
  
  \tagitem{prvTrue}{$\Gamma \Proves[\Sigma] \ltrue$, अतः निगमनात्मक
    संवृत्ति, अर्थात् स्थिति~\olref{prop:ccs-closed}, से $\ltrue \in
    \Gamma$।}{}

\item यदि $\lnot!A \in \Gamma$, तो संगति से $!A \notin \Gamma$; और यदि
    $!A \notin \Gamma$, तो $!A \in \Gamma$, क्योंकि $\Gamma$ पूर्ण
    $\Sigma$-संगत है।
    
  \tagitem{prvAnd}{\iftag{probAnd}{अभ्यास।}{मान लें $!A \land !B \in
      \Gamma$। चूँकि $(!A \land !B) \lif !A$ सर्वसत्यतात्मक दृष्टांत है,
      निगमनात्मक संवृत्ति, अर्थात् स्थिति~\olref{prop:ccs-closed}, से
      $!A \in \Gamma$। इसी प्रकार $!B \in \Gamma$। दूसरी ओर, मान लें कि
      $!A \in \Gamma$ और $!B \in \Gamma$, दोनों। तब निगमनात्मक संवृत्ति
      से $(!A \land !B) \in \Gamma$, क्योंकि $!A \lif (!B \lif (!A
      \land !B))$ सर्वसत्यतात्मक दृष्टांत है।}}{}

  \tagitem{prvOr}{\iftag{probOr}{अभ्यास।}{मान लें $!A \lor !B \in
      \Gamma$, तथा $!A \notin \Gamma$ और $!B \notin \Gamma$। चूँकि
      $\Gamma$ पूर्ण $\Sigma$-संगत है, $\lnot!A \in \Gamma$ और $\lnot
      !B \in \Gamma$। तब $\lnot(!A \lor !B) \in \Gamma$, क्योंकि
      $\lnot !A \lif (\lnot !B \lif \lnot (!A \lor !B))$
      सर्वसत्यतात्मक दृष्टांत है। इसका अर्थ होगा कि $\Gamma$,
      $\Sigma$-असंगत है, जो विरोधाभास है।}}{}
  
  \tagitem{prvIf}{\iftag{probIf}{अभ्यास।}{मान लें $!A \lif !B \in
      \Gamma$ और $!A \in \Gamma$; तब $\Gamma \Proves[\Sigma] !B$, और
      इसलिए निगमनात्मक संवृत्ति से $!B \in \Gamma$। विलोमतः, यदि $!A
      \lif !B \notin \Gamma$, तो $\Gamma$ के पूर्ण $\Sigma$-संगत होने
      से $\lnot (!A \lif !B) \in \Gamma$। चूँकि $\lnot(!A \lif !B)
      \lif !A$ सर्वसत्यतात्मक दृष्टांत है, निगमनात्मक संवृत्ति से $!A
      \in \Gamma$। चूँकि $\lnot(!A \lif !B) \lif \lnot !B$
      सर्वसत्यतात्मक दृष्टांत है, $\lnot !B \in \Gamma$। तब $\Gamma$ के
      $\Sigma$-संगत होने से $!B \notin \Gamma$।}}{}
  
  \tagitem{prvIff}{\iftag{probIff}{अभ्यास।}{मान लें $!A \liff !B \in
      \Gamma$। यदि $!A \in \Gamma$, तो $!B \in \Gamma$, क्योंकि $(!A
      \liff !B) \lif (!A \lif !B)$ सर्वसत्यतात्मक दृष्टांत है। इसी प्रकार,
      यदि $!B \in \Gamma$, तो $!A \in \Gamma$। अतः या तो $!A \in
      \Gamma$ और $!B \in \Gamma$, दोनों, अथवा न $!A \in \Gamma$ और न
      $!B \in \Gamma$।
      
      विलोमतः, मान लें $!A \lif !B \notin \Gamma$। चूँकि $\Gamma$ पूर्ण
      $\Sigma$-संगत है, $\lnot (!A \liff !B) \in \Gamma$। चूँकि
      $\lnot(!A \liff !B) \lif (!A \lif \lnot !B)$ सर्वसत्यतात्मक
      दृष्टांत है, यदि $!A \in \Gamma$, तो $\lnot !B \in \Gamma$; और
      $\Gamma$ के $\Sigma$-संगत होने से $!B \notin \Gamma$। इसी प्रकार,
      यदि $!B \in \Gamma$, तो $!A \notin \Gamma$। अतः न तो $!A \in
      \Gamma$ और $!B \in \Gamma$, और न ही $!A \notin \Gamma$ और $!B
      \notin \Gamma$।}}{}
  \end{enumerate}
\end{proof}

\begin{probtag}{probAnd,probOr,probIf,probIff}
  \olref[nml][com][ccs]{prop:ccs-properties} का प्रमाण पूरा कीजिए।
\end{probtag}

\end{document}
